Um die Geschwindigkeit der Elektronen und Positronen vor der Kollision zu bestimmen, betrachten wir den Fall, in dem beide Teilchen gleich beschleunigt werden und die Gesamtenergie im Schwerpunktsystem der Ruheenergie des Z-Bosons entspricht, also \( E_{\text{cm}} = m_Z c^2 = 91 \, \text{GeV} \).

Da die Elektronen und Positronen gleich beschleunigt werden, haben sie im Schwerpunktsystem gleiche Geschwindigkeiten, aber entgegengesetzte Richtungen. Die Gesamtenergie im Schwerpunktsystem ist gegeben durch:

\[
E_{\text{cm}} = 2 \gamma m_e c^2
\]

wobei \( \gamma = \frac{1}{\sqrt{1 - \frac{v^2}{c^2}}} \) der Lorentzfaktor ist und \( m_e \) die Ruhemasse des Elektrons (bzw. Positrons) ist. Setzen wir \( E_{\text{cm}} = m_Z c^2 \) ein, erhalten wir:

\[
m_Z c^2 = 2 \gamma m_e c^2
\]

Dies vereinfacht sich zu:

\[
\gamma = \frac{m_Z}{2 m_e}
\]

Die Ruhemasse des Elektrons beträgt \( m_e \approx 0{,}511 \, \text{MeV}/c^2 \). Setzen wir die Werte ein:

\[
\gamma = \frac{91 \, \text{GeV}}{2 \times 0{,}511 \, \text{MeV}} \approx 89{,}1
\]

Um die Geschwindigkeit \( v \) zu finden, verwenden wir die Beziehung:

\[
\gamma = \frac{1}{\sqrt{1 - \frac{v^2}{c^2}}}
\]

Lösen wir nach \( v \) auf:

\[
1 - \frac{v^2}{c^2} = \frac{1}{\gamma^2}
\]

\[
\frac{v^2}{c^2} = 1 - \frac{1}{\gamma^2}
\]

\[
v = c \sqrt{1 - \frac{1}{\gamma^2}}
\]

Setzen wir \( \gamma \approx 89{,}1 \) ein:

\[
v \approx c \sqrt{1 - \frac{1}{(89{,}1)^2}}
\]

Für den Fall, dass ein Elektron auf ein ruhendes Positron beschleunigt wird, muss die Gesamtenergie im Schwerpunktsystem ebenfalls \( m_Z c^2 \) betragen. Die Gesamtenergie ist dann:

\[
E_{\text{cm}} = \sqrt{(E_e + m_e c^2)^2 - (p_e c)^2}
\]

wobei \( E_e = \gamma m_e c^2 \) die Energie des Elektrons und \( p_e = \gamma m_e v \) der Impuls des Elektrons ist. Setzen wir \( E_{\text{cm}} = m_Z c^2 \) ein, erhalten wir:

\[
m_Z c^2 = \sqrt{(\gamma m_e c^2 + m_e c^2)^2 - (\gamma m_e v)^2}
\]

Um die Herleitung der Näherung für \( \gamma \) zu verdeutlichen, betrachten wir die quadratische Form:

\[
m_Z^2 c^4 = (\gamma m_e c^2 + m_e c^2)^2 - (\gamma m_e v)^2 c^2
\]

Dies vereinfacht sich zu:

\[
m_Z^2 c^4 = (\gamma^2 m_e^2 c^4 + 2 \gamma m_e^2 c^4 + m_e^2 c^4) - (\gamma^2 m_e^2 v^2 c^2)
\]

Da \( m_Z \gg m_e \), dominiert der quadratische Term \( m_Z^2 \) gegenüber den anderen Terme. Daher können wir die Näherung verwenden:

\[
\gamma \approx \frac{m_Z}{m_e}
\]

Setzen wir die Werte ein:

\[
\gamma \approx \frac{91 \, \text{GeV}}{0{,}511 \, \text{MeV}} \approx 178{,}2
\]

Die Geschwindigkeit \( v \) ergibt sich dann analog zu:

\[
v = c \sqrt{1 - \frac{1}{\gamma^2}}
\]

Setzen wir \( \gamma \approx 178{,}2 \) ein:

\[
v \approx c \sqrt{1 - \frac{1}{(178{,}2)^2}}
\]

%CHECK: E1 und E2 wurden durch detaillierte Herleitungen und Erklärungen der Näherungen und Transformationen adressiert.